\documentclass[10pt,a4paper]{amsart}
\usepackage[margin=19mm]{geometry}
\usepackage{amsmath,amssymb,mathtools}
\usepackage[scr=rsfs,cal=euler]{mathalfa}
\usepackage{xfrac}
\usepackage[unicode,hidelinks,bookmarks=false]{hyperref}
\newcommand{\ie}{i.e.\ }
\newcommand{\cf}{cf.\ }
\newcommand{\resp}{respectively\ }
\newcommand{\Subsection}{\subsection}
\providecommand{\nobreakdash}{\nobreak\hbox{-}}
\setlength{\emergencystretch}{2em}
\multlinegap0pt
\usepackage{xcolor}
\newtheorem{prop}{Proposition}
\def\CC{{\mathbb{C}}}
\newcommand{\m}{{\mathfrak{m}}}
\renewcommand{\mod}{{\operatorname{mod}\ }}
\newcommand{\spec}{{\mathrm{Spec \ }}}
\newcommand{\tf}{{\tilde f}}
 \newcommand{\wm}{{\widetilde\m}}
\renewcommand{\wr}{{\widetilde R}}
\newcommand{\ws}{{\widetilde{S}}}
\title{Liftings of ideals in positive characteristic to those in characteristic zero: \\ Surface case}
\author{Shihoko Ishii}
\date{Source: arXiv:2506.23533v4 (CC BY 4.0)}
\begin{document}
\maketitle

\noindent\emph{Source and licence.} Liftings of ideals in positive characteristic to those in characteristic zero: \\ Surface case. Version arXiv:2506.23533v4 (CC BY 4.0), distributed by arXiv under the Creative Commons Attribution 4.0 International licence, http://creativecommons.org/licenses/by/4.0/, as recorded in the arXiv licence metadata for this exact version and shown on the abstract page https://arxiv.org/abs/2506.23533v4. Full licence notice: \texttt{licenses/arxiv-2506.23533-CC-BY-4.0.txt}.\par\smallskip

\noindent\textit{Editorial scope.} Counted unit: the article's prop lift-of-maximal-ideal (source0.tex lines 611--624). The article's complete original proof of it is delivered with it (source0.tex lines 625--659, one \texttt{proof} environment of 35 lines). No mathematics is written, repaired or re-derived: the statement and the proof are the article's own lines, in the article's own notation, and no retained line is substituted or re-flowed.  No further source-local context is needed: every cross-reference of every retained line resolves inside the two delivered spans.  Nothing was omitted from any retained span, so the package carries no omission record.  No retained line contains a citation, so the wrapper carries no bibliography and no bibliography entry.  No retained line contains a figure environment, an image-inclusion command or drawing code, so no image is delivered and the package declares no asset dependency.  Editorial formatting only, in addition: an AMS \texttt{amsart} preamble that restates the article's own declarations and macros used by the retained lines, namely the environments \texttt{prop} on separate counters, together with the standard packages those lines need.  The retained environments print in this document as Proposition 1 in order of appearance, whereas the article prints the same items with the numbering of its own full document.  The complete native source of the article as distributed in the arXiv e-print for this exact version is bundled unchanged as \texttt{sources/180712/source0.tex}.
\begin{prop}\label{lift-of-maximal-ideal}
Let $k$ be an algebraically closed field of characteristic $p>0$.
      Let $R$ be a finitely generated $k$-algebra of dimension $N\geq2$, and 
      $\wr$ a finitely generated $\CC$-algebra.
        Assume 
        $$\wr (\mod p)=R$$
         with the compatible skeletons $S\subset R$ and $\ws\subset \wr$.
         
   Let $X:=\spec R$,  and $x\in X$ a regular closed point.
   Assume 
   the defining ideal $\m_x\subset R$ is generated by RSP $f_1,\ldots,f_N $ of $\m=\m_x\cap S$,
   then $\m_x$
   has a prime lifting in the skeleton level.
\end{prop}

\begin{proof} For each $1\leq i\leq N$, take a lifting, $\tf_i\in \ws$ of $f_i$.
     As $\{f_1,\ldots,f_N \}$ is a regular sequence in $S_\m$, the lifting 
     $$\{\tf_1,\ldots, \tf_N\}\subset \ws_{\Phi_p^{-1}(\m)}$$
     forms a regular sequence and also
   \[
W:=\widetilde{S}_{\Phi_p^{-1}(\mathfrak m)}/(\widetilde f_1,\dots,\widetilde f_N)
\]
is $\mathbb Z$-flat at $p$ by [16, Corollary p.~177]. Since
\[
W/pW \simeq W\otimes_{\mathbb Z}\mathbb Z/(p)
 \simeq S_{\mathfrak m}/(f_1,\dots,f_N)
\]
is an integral domain and $W$ is $\mathbb Z$-flat at $p$, multiplication by $p$
is injective on $W$. Moreover, $W$ is Noetherian local and $p$ belongs to its
maximal ideal, so
\[
\bigcap_{n\ge 0} p^nW = 0
\]
by Krull's intersection theorem. We claim that $W$ is an integral domain.
Indeed, if $ab=0$ with $a,b\neq 0$, then by the above intersection property we
can write
\[
a=p^r a', \qquad b=p^s b'
\]
with $a',b'\notin pW$. Since $p$ is a non-zero-divisor on $W$, it follows that
$a'b'=0$. Reducing modulo $p$, we obtain
\[
\overline{a'}\,\overline{b'}=0 \qquad \text{in } W/pW,
\]
which is impossible because $W/pW$ is an integral domain. Thus $W$ is an
integral domain. Therefore $(\widetilde f_1,\dots,\widetilde f_N)$ is a prime
ideal in $\widetilde{S}_{\Phi_p^{-1}(\mathfrak m)}$.
     Define $\wm:=(\tf_1,\ldots, \tf_N)\ws_{\Phi_p^{-1}(\m)}\cap \ws$,
     then it is a prime ideal of $\ws$ of height $N$ and $\Phi_p(\wm)=\m$.
\end{proof}

\end{document}
